\chapter{Rejected Design: On-Device Models}
\label{chap:mobile_llms}

\section{What Was Considered}

Early in the project, we considered running a small quantised model on the
same machine as the Lean server, eliminating the network entirely.
The leading candidate was \textbf{Qwen2.5-Math-1.5B-Instruct} in GGUF Q4
format ($\approx$1\,GB), runnable offline via \texttt{llama.cpp} or
\texttt{MLC-LLM}.

The reasoning was sound: no API outages, no rate limits, fully reproducible,
no cost. The Lean REPL is local anyway, so the only inter-process hop would
be a loopback socket.

\section{Why We Did Not Use It}

\begin{enumerate}
  \item \textbf{Capability gap.} The research question is what a
        \emph{frontier} model achieves on real Mathlib theorems.
        A 1.5B parameter model cannot answer that question; it can only
        establish a floor.
  \item \textbf{Cache dependency.} The Lean build cache must be warm for
        every run --- 8{,}667 \texttt{.olean} files, and a \texttt{.lake}
        directory of roughly 6.5\,GB across 140{,}000 files once dependencies
        are unpacked. A machine provisioned for that, and for a resident Lean
        REPL holding a full Mathlib environment in memory, is not meaningfully
        constrained by a 1\,GB model weight. The on-device design's central
        advantage --- a small footprint --- buys nothing here.
  \item \textbf{Capability was the deciding factor, not reproducibility.}
        We chose hosted APIs for reach, accepting a reproducibility cost
        rather than avoiding one. See below.
\end{enumerate}

\section{The Reproducibility Cost We Accepted}

It would be convenient to claim that versioned API endpoints give the same
reproducibility guarantee as local weights. They do not, and this project has
direct evidence against it.

\begin{itemize}
  \item \textbf{Models were withdrawn mid-study.} A smoke test against
        \texttt{llama-3.1-8b-instant} returned \texttt{404 model\_not\_found};
        \texttt{llama3-8b-8192} and \texttt{mixtral-8x7b-32768} returned
        \texttt{400 model\_decommissioned}. Model identifiers that were
        current when the harness was written had ceased to exist by the time
        it ran against them.
  \item \textbf{Endpoints went down for days, not minutes.} The NVIDIA NIM
        endpoint began returning 404/429/503 partway through an Easy-50 run
        and had not recovered three days later, costing 24 of 50 theorems.
  \item \textbf{Provider-side behaviour changed under us.} Groq's harmony-format
        models intermittently returned \texttt{400 tool\_use\_failed} --- the
        provider's own parser misreading its own model's output --- which the
        harness now has to detect and recover from.
\end{itemize}

An offline GGUF checkpoint has none of these properties: the weights are a
file, and a file does not get deprecated, rate-limited, or rewritten upstream.
The honest statement of the trade is therefore that we gave up
reproducibility and uptime in exchange for evaluating a model 350 times
larger, and that the cost was real --- roughly half the theorems in each of
the two Easy-50 runs were lost to provider problems rather than to the model
failing at mathematics.

This is why every result in this study reports a reachable-only pass rate
alongside the headline: the headline figure silently mixes provider uptime
into a capability measurement, and the gap between the two numbers is exactly
the price of this design decision.

\section{Future Use}

A small on-device model remains a natural ablation: does model scale predict
proof rate, or does the failure taxonomy look the same at 1.5B and 550B?
This is deferred to Phase~2 (multi-model comparison).

The instrumentation is already in place. Provider selection is a runtime flag
rather than a code change, and any OpenAI-compatible endpoint works --- which
includes \texttt{llama.cpp}'s own server. Adding a local model is a matter of
pointing the base URL at \texttt{localhost} and setting an appropriate
completion budget, not of modifying the harness.
